\documentclass[11pt]{article}
\usepackage[margin=1in]{geometry}
\usepackage{amsmath,amssymb}
\usepackage{booktabs}
\usepackage[T1]{fontenc}
\usepackage{lmodern}

\title{Zero Day -- Test 242: Numerical Robustness of the Bounce and Recollapse}
\author{Inacio F. Vasquez\\Independent Researcher, United States}
\date{20 October 2026}

\begin{document}
\maketitle

\section{Scope}
This note records a supplied three-run solver comparison for an illustrative
zero-charge, closed-FLRW trajectory. The integration data and solver outputs
have not been independently reproduced as part of this document. The governing
equations are assumed for this calculation and have not here been derived from
the DFM--MKC action.

\section{Independent-method comparison}
\begin{center}
\begin{tabular}{lrrr}
\toprule
Solver & Tolerance & Maximum-size time & Time to $a=0.1$\\
\midrule
DOP853 & $10^{-10}$ & 397.749694 & 792.771210\\
DOP853 & $10^{-12}$ & 397.749690 & 792.771146\\
RK45   & $10^{-10}$ & 397.750428 & 792.782148\\
\bottomrule
\end{tabular}
\end{center}

The maximum-size time spans
\[
397.750428-397.749690=0.000738
\]
across the three reported runs. The time to $a=0.1$ spans
\[
792.782148-792.771146=0.011002.
\]
These ranges describe agreement among the reported configurations; they are
not rigorous error bounds and do not by themselves establish convergence of
the trajectory or its endpoint.

\section{Result}
\textbf{RESULT:} The maximum-size turnaround time is numerically stable across
the reported solver settings. The late contracting phase and the endpoint as
$a\to0$ remain unresolved.

Adaptive integration tolerances control an estimated numerical error for the
integrator. They do not validate the physical equations, the potential, or the
model assumptions.

\section{Weakest missing result}
A controlled endpoint study is needed. Monitor curvature invariants, including
the Kretschmann scalar
\[
K=R_{\mu\nu\rho\sigma}R^{\mu\nu\rho\sigma},
\]
alongside $a(t)$ and $H(t)$. Repeat the contracting-phase integration with
appropriate independent solvers and endpoint-aware variables, and demonstrate
convergence over a sequence of smaller scale factors. A divergence of $K$ as
$a\to0$, if established under the stated equations, would support curvature
singularity classification; it would not alone prove every form of geodesic
incompleteness. If $K$ remains bounded, further invariants and extension
analysis may still be required.

\section{Boundary}
\[
\texttt{BOUNDARY :=}
\]
This calculation does not derive its closed-FLRW equations from the DFM--MKC
action, does not constitute an empirical cosmology fit, and does not prove a
beginningless universe or geodesic incompleteness.

\section{Next actions}
\begin{enumerate}
  \item Preserve the raw integration configuration, parameter values, and output.
  \item Check curvature invariants along the contracting branch.
  \item Test endpoint convergence with independent numerical methods.
  \item Derive the governing equations from the action before extending the claim.
\end{enumerate}

\end{document}
